% Fragmento público generado desde propietarios íntegros.
% Residencia: 52.13.1.
% Fuente: ../incoming/radion_monografia_20260827/manuscrito/sections/13_sintesis.tex:81-134
\subsubsection{Définition opératoire du radion}

\begin{definition}[Radion HMT--MD]
Le radion est la section orientée de phase, d'action et de mémoire qui satisfait
simultanément :
\begin{enumerate}
\item parcourt une unité de phase \(\theta=1\) ;
\item balaie une action signée \(\sigma\hret\) ;
\item appartient à un tour \(T_+=2\pih\) généré coinductivement ;
\item conserve le feuillet \(\sigma\), le retour nonadique et le cylindre de
raffinement ;
\item satisfait le raccord exact
\(1=S_E-\Phi_T+\epsone\).
\end{enumerate}
\end{definition}

Le radian conventionnel est l'image du radion par le foncteur d'oubli
\[
\mathfrak R_\sigma
\longmapsto
\theta=1\in\R/2\pi\Z.
\]
Cet oubli est utile et légitime pour la trigonométrie euclidienne. Il devient
insuffisant lorsque la physique doit distinguer action, signe, retour,
champ et histoire.

\paragraph{Généalogie de l'unité angulaire naturelle}

Le caractère naturel du radian est généralement justifié par le fait que la dérivée de
\(e^{i\theta}\), \(\sin\theta\) et \(\cos\theta\) adopte sa forme la plus simple
lorsque \(\theta\) est mesuré comme quotient arc/rayon. HMT conserve cet avantage,
mais l'explique depuis une couche antérieure :
\[
\theta\quad\leftrightarrow\quad
\frac{\mathcal A(\theta)}{\hret}.
\]
La phase est naturelle parce qu'elle mesure l'action en unités de \(\hret\) ; le tour est
naturel parce qu'il enferme \(\hfull\). Le quotient arc/rayon est la publication
euclidienne d'une unité qui possède déjà une structure symplectique et une mémoire.

\paragraph{Compatibilité entre ellipse d'action et croissance dimensionnelle}

Le cercle est la forme qui résulte de l'oubli de l'anisotropie :
\(C^*=0\Rightarrow\eta=0\Rightarrow a=b\). L'ellipse est plus informative parce qu'elle
conserve deux valeurs propres, deux axes et une orientation d'échange à produit
fixe. En ce sens, l'ellipse est une prolongation du cercle, non un cercle
déformé par accident.

La profondeur de Klein et le relèvement hélicoïdal ajoutent deux types d'information
que le contour plan ne contient pas : distance intérieure et mémoire de retour.
Telle est l'image ontologique de la mécanique dimensionnelle qui émerge du radion :
une dimension nouvelle apparaît lorsque l'on exige de conserver une donnée indépendante
que la projection précédente devait oublier.

